\documentclass{article}
\usepackage[margin=1in]{geometry}
\usepackage{amsfonts}
\usepackage{amsmath}
\usepackage{amsthm}
\usepackage{parskip}
\usepackage{enumitem}
\usepackage{mathtools}
\usepackage{booktabs}
\usepackage{amssymb}
\newtheorem{theorem}{Theorem}

\DeclarePairedDelimiterX{\norm}[1]{\lVert}{\rVert}{#1}
\newcommand{\Omh}{\hat{\Omega}}
\newcommand{\Ups}{\Upsilon}

\title{Factored Low-Rank Decomposition of $\exp(\Phi)$\\
\large Connection to the Normal Operator Toeplitz Method}
\author{viktor.vigren}
\date{\today}

\begin{document}
\maketitle

This document complements \texttt{MRI\_Physics.tex} by addressing a key difficulty with the
joint SVD approach: the effective rank of $\exp(\Phi)$ grows multiplicatively when off-resonance
and nonlinear gradient fields are combined, forcing $L \gg 1$ for accurate reconstruction.
We show that a \emph{factored} SVD decomposition achieves the same accuracy with $L$ that grows
\emph{additively} and that the resulting compound temporal basis slots directly into the Toeplitz
normal-operator framework.

% ---------------------------------------------------------------------------
\section{Readout-Relative Formulation}

The full phase modulation from \texttt{MRI\_Physics.tex} is
\begin{equation}
    \Phi(\mathbf{r}, t) = z(\mathbf{r})\,t +
        \sum_q b_q(\mathbf{r})\,\alpha_q(t),
    \qquad
    z(\mathbf{r}) = -\tfrac{1}{T_2(\mathbf{r})} - i\,\gamma\,\Delta B_0(\mathbf{r}),
\end{equation}
where $\alpha_q(t) = \int_0^t G_q(\tau)\,d\tau$ is the $q$-th nonlinear gradient
waveform integrated from the start of the readout.

\paragraph{Pre-readout phases are irrelevant.}
Any field contribution $G_q(\tau)$ applied \emph{before} readout (e.g.\ velocity-encoding
lobes, dephasing gradients) integrates to a value $\alpha_q^{\rm pre}$ that is
\emph{identical for every readout sample within a given spoke}.  It therefore
multiplies $\exp(\Phi)$ by a constant $e^{b_q(\mathbf{r})\,\alpha_q^{\rm pre}}$
which is absorbed into the (unknown) image $\rho(\mathbf{r})$ and plays no role in
the forward model used for reconstruction.
Consequently $t$ in the above is measured from the \emph{start of the readout
window only} and $\alpha_q(t)$ includes only the readout-gradient contribution.

For a koosh-ball (3D radial) trajectory the readout gradient is applied along spoke
direction $\hat{d}_s$, so the dominant nonlinear waveforms (Maxwell concomitant
fields up to second order) factorize as
\begin{equation}
    \alpha_q(k{=}(s,j)) = d_q^{(s)}\, f_q(j),
    \label{eq:alpha_factor}
\end{equation}
where $d_q^{(s)}$ is a scalar that depends only on spoke direction (e.g.\
$d_{z^2}^{(s)} = \cos^2\theta_s$) and $f_q(j)$ is the waveform shape shared by
all spokes (e.g.\ a bipolar profile in readout index $j$).

% ---------------------------------------------------------------------------
\section{Why Joint SVD Requires Large $L$}

The joint SVD seeks a rank-$L$ factorization
$\exp(\Phi(\mathbf{r}_j, t_k)) \approx \sum_{l=1}^L \Omega_{kl}\,\Upsilon_{jl}$.
The effective rank of $\exp(\Phi)$ is bounded by the product of the effective ranks
of its factors:
\begin{equation}
    \operatorname{rank}^*\bigl(\exp(\Phi)\bigr)
    \;\leq\; \operatorname{rank}^*\bigl(\exp(z\,t)\bigr) \;\times\;
    \prod_q \operatorname{rank}^*\bigl(\exp(i\,b_q\,\alpha_q)\bigr).
\end{equation}
Here $\operatorname{rank}^*(\cdot)$ denotes the number of singular values needed
to achieve a target relative error $\varepsilon$.
Fessler's result (Sutton et al.\ 2003) shows that $\operatorname{rank}^*(\exp(-i\,2\pi\,\Delta B_0\,t)) \approx L_o \sim 7$
for a $140\,\text{Hz}$ B0 range and $18\,\text{ms}$ readout.
Each nonlinear factor with max phase $\pi$ requires $L_q \sim 4$ additional terms
(Chebyshev/PSWF argument: rank scales as $2\Delta\phi/\pi + 1$).

With $Q=2$ nonlinear fields the joint SVD therefore requires
$L \approx L_o \times L_1 \times L_2 = 7 \times 4 \times 4 = 112$ for
$\varepsilon = 10^{-4}$, making iterative reconstruction with $20 \times 8 \times 5 \times 32 \times 112$
Fourier transforms per iteration infeasible.

% ---------------------------------------------------------------------------
\section{Factored SVD Decomposition}

\subsection{Decomposing $\exp(\Phi)$ into independent factors}

Write
\begin{equation}
    \exp\bigl(\Phi(\mathbf{r}, t_k)\bigr)
    = \underbrace{e^{z(\mathbf{r})\,t_k}}_{\displaystyle A_o(\mathbf{r},k)}
    \;\prod_{q=1}^Q
    \underbrace{e^{i\,b_q(\mathbf{r})\,\alpha_q(k)}}_{\displaystyle A_q(\mathbf{r},k)}.
    \label{eq:phi_factored}
\end{equation}

Each factor involves \emph{only one field type} and its corresponding waveform.

\subsection{Individual histogram and SVD per factor}

For each factor we define a separate histogram over the relevant feature.

\paragraph{Off-resonance factor $A_o$.}
Feature: $z(\mathbf{r})$ only.  Histogram bins $h_o \in \{1,\ldots,n_o\}$ with
representative values $\bar{z}_{h_o}$.  The $K \times n_o$ matrix
$[\Phi_o]_{k,h_o} = e^{\bar{z}_{h_o}\,t_k}$
has SVD-optimal rank $L_o \sim 7$ for typical parameters.
\begin{equation}
    e^{\bar{z}_{h_o}\,t_k} \;\approx\; \sum_{l_o=1}^{L_o} \Omega_o[k,\,l_o]\;\Upsilon_o[h_o,\,l_o].
    \label{eq:svd_offres}
\end{equation}

\paragraph{Nonlinear factor $A_q$.}
Feature: $b_q(\mathbf{r})$ only.  Histogram bins $h_q \in \{1,\ldots,n_q\}$ with
representative values $\bar{b}_q^{(h_q)}$.  The $K \times n_q$ matrix
$[\Phi_q]_{k,h_q} = e^{i\,\bar{b}_q^{(h_q)}\,\alpha_q(k)}$
has rank $L_q \sim 4$ for max phase $\leq \pi$.
\begin{equation}
    e^{i\,\bar{b}_q^{(h_q)}\,\alpha_q(k)} \;\approx\; \sum_{l_q=1}^{L_q} \Omega_q[k,\,l_q]\;\Upsilon_q[h_q,\,l_q].
    \label{eq:svd_nl}
\end{equation}

Note that the histograms are now \emph{one-dimensional per factor}: $n_o = n_z^2$
(complex $z$ values) and $n_q = n_{b_q}$ (real field values).  This is far smaller
than the joint histogram $n_{\rm joint} = n_z^2 \times \prod_q n_{b_q}^2$.

\subsection{Combined forward operator}

Substituting \eqref{eq:svd_offres}--\eqref{eq:svd_nl} into \eqref{eq:phi_factored}:
\begin{equation}
    \exp\bigl(\Phi(\mathbf{r}_j, t_k)\bigr) \;\approx\;
    \sum_{\boldsymbol{l}} \Omega_{\boldsymbol{l}}(k)\;\Upsilon_{\boldsymbol{l}}(j),
    \label{eq:factored_approx}
\end{equation}
where the multi-index $\boldsymbol{l} = (l_o, l_1, \ldots, l_Q)$ runs over
$L_{\rm fwd} = L_o \times \prod_q L_q$ compound terms and
\begin{align}
    \Omega_{\boldsymbol{l}}(k)   &= \Omega_o[k,l_o]\;\prod_q \Omega_q[k,l_q], \\
    \Upsilon_{\boldsymbol{l}}(j) &= \Upsilon_o[\mathrm{bin}_o(j),\,l_o]\;\prod_q \Upsilon_q[\mathrm{bin}_q(j),\,l_q].
\end{align}

The forward signal is then
\begin{equation}
    S_c(k) \approx \sum_{\boldsymbol{l}}
    \Omega_{\boldsymbol{l}}(k)\;
    \mathcal{F}\!\left[\,\rho(\mathbf{r})\,s_c(\mathbf{r})\,\Upsilon_{\boldsymbol{l}}(\mathbf{r})\,\right]\!(k),
    \label{eq:forward_factored}
\end{equation}
requiring $L_{\rm fwd} = L_o \times \prod_q L_q$ NUFFT evaluations.

% ---------------------------------------------------------------------------
\section{Error Analysis: Additive Bound for Unitary Factors}

\begin{theorem}[Additive error for factored unitary approximation]
Let $A_q(\mathbf{r},k)$ be the exact factors and $\hat{A}_q(\mathbf{r},k)$ their
rank-$L_q$ approximations with errors
$\varepsilon_q = \norm{A_q - \hat{A}_q}_F / \norm{A_q}_F$.
Since $|A_q(\mathbf{r},k)| = 1$ for all $(\mathbf{r},k)$
(purely imaginary phases and pure T2 decay satisfying $|e^{z\,t}|\leq 1$),
\begin{equation}
    \norm{\exp(\Phi) - \hat{A}_o \prod_q \hat{A}_q}_F
    \;\leq\;
    \varepsilon_o \norm{\exp(\Phi)}_F \;+\; \sum_q \varepsilon_q \norm{\exp(\Phi)}_F.
    \label{eq:additive_bound}
\end{equation}
\end{theorem}

\begin{proof}
By telescoping.  Write
$A_o \prod_q A_q - \hat{A}_o \prod_q \hat{A}_q$
as a telescoping sum of replacements $A_q \to \hat{A}_q$ one factor at a time.
Each replacement contributes a term of the form
$(\hat{A}_1 \cdots \hat{A}_{q-1})(A_q - \hat{A}_q)(A_{q+1} \cdots A_Q)$
whose Frobenius norm is bounded by $\varepsilon_q \norm{\exp(\Phi)}_F$
using $|A_p| = 1$ and $|\hat{A}_p| \leq 1 + \varepsilon_p \approx 1$ for all other factors $p$.
Summing over all replacements gives \eqref{eq:additive_bound}.
\end{proof}

\paragraph{Contrast with joint SVD.}
The joint SVD minimises $\norm{\exp(\Phi) - \hat{\Phi}_{\rm joint}}_F$ directly, but
the effective rank of the product is \emph{at least} $\operatorname{rank}^*(A_o) \times \prod_q \operatorname{rank}^*(A_q)$
because the singular vectors of $A_o \odot A_q$ are not the tensor products of the
individual singular vectors.  In practice the joint SVD at $L=30$ achieves the same
Frobenius error as the factored approach at $L_o + L_1 + L_2 = 15$, but the joint
SVD must spend all $L$ vectors simultaneously capturing both off-resonance and
nonlinear dynamics.  In the factored approach, off-resonance always gets its full
$L_o = 7$ dedicated vectors, guaranteeing Fessler-level accuracy \emph{independently}
of the nonlinear correction.

% ---------------------------------------------------------------------------
\section{Normal Operator with Factored Compound Basis}

\subsection{Naive compound normal operator}

From \eqref{eq:factored_approx}, the normal operator inherits
\begin{equation}
    \exp(\Phi_j + \Phi_i^*) \approx
    \sum_{\boldsymbol{l}, \boldsymbol{m}}
    \Omega_{\boldsymbol{l}}(k)\,\Omega_{\boldsymbol{m}}^*(k)\;
    \Upsilon_{\boldsymbol{l}}(j)\,\Upsilon_{\boldsymbol{m}}^*(i),
\end{equation}
with $L_{\rm fwd}^2 = (L_o \times \prod_q L_q)^2$ compound pairs.

\subsection{Exploiting the product structure in the secondary compression}

Because $\Omega_{\boldsymbol{l}}(k) = \prod_{f} \Omega_f[k, l_f]$ (where $f$ indexes
all factors including $o$ and each $q$), the product of two compound temporal
coefficients also factorizes:
\begin{equation}
    \Omega_{\boldsymbol{l}}(k)\,\Omega_{\boldsymbol{m}}^*(k) =
    \prod_{f} \Bigl[\Omega_f[k,l_f]\,\Omega_f^*[k,m_f]\Bigr].
    \label{eq:omega_product}
\end{equation}

We can therefore apply the secondary compression of \texttt{MRI\_Physics.tex}
\emph{factor by factor}.  For each factor $f$, form the $L_f(L_f+1)/2 \times K$
matrix of upper-triangular products and compress to $P_f$ terms:
\begin{equation}
    \Omega_f[k,l_f]\,\Omega_f^*[k,m_f] \;\approx\;
    \sum_{p_f=1}^{P_f} C^{(f)}_{l_f m_f,\,p_f}\;\Lambda^{(f)}_{k,\,p_f}.
    \label{eq:factor_svd2}
\end{equation}

Substituting into \eqref{eq:omega_product}:
\begin{equation}
    \Omega_{\boldsymbol{l}}(k)\,\Omega_{\boldsymbol{m}}^*(k) \approx
    \sum_{\boldsymbol{p}}
    \underbrace{\prod_f C^{(f)}_{l_f m_f,\,p_f}}_{\displaystyle C_{\boldsymbol{l}\boldsymbol{m},\boldsymbol{p}}}
    \underbrace{\prod_f \Lambda^{(f)}_{k,\,p_f}}_{\displaystyle \Lambda_{k,\boldsymbol{p}}},
\end{equation}
where the multi-index $\boldsymbol{p} = (p_o, p_1, \ldots, p_Q)$ runs over
$P_{\rm total} = \prod_f P_f$ terms.

\subsection{Rank decomposition of $C$}

Each coefficient matrix $C_{\boldsymbol{l}\boldsymbol{m},\boldsymbol{p}}$ is a rank-$R$
approximation following the same structure as \texttt{MRI\_Physics.tex}:
\begin{equation}
    C_{\boldsymbol{l}\boldsymbol{m},\boldsymbol{p}} \approx \sum_{r=1}^R
    \hat{C}_{\boldsymbol{l},r,\boldsymbol{p}}\;\hat{C}^*_{\boldsymbol{m},r,\boldsymbol{p}}.
\end{equation}

However, because $C$ inherits the product structure \eqref{eq:factor_svd2}, we can
decompose \emph{each factor's} $C^{(f)}$ matrix with rank $R_f$, giving
\begin{equation}
    \hat{C}_{\boldsymbol{l},r,\boldsymbol{p}} = \prod_f \hat{C}^{(f)}_{l_f,\,r_f,\,p_f},
    \quad R_{\rm total} = \prod_f R_f.
\end{equation}

\subsection{Final normal operator}

Collecting all terms (parallel to eq.~(6) of \texttt{MRI\_Physics.tex}):
\begin{align}
    \exp(\Phi_j + \Phi_i^*) &\approx
    \sum_{\boldsymbol{p},r} \Lambda_{k,\boldsymbol{p}}\;\Theta_{j,r,\boldsymbol{p}}\;\Theta^*_{i,r,\boldsymbol{p}},
    \label{eq:final_normal_factored}
\end{align}
where the combined spatial basis vectors are
\begin{equation}
    \Theta_{j,r,\boldsymbol{p}} = \sum_{\boldsymbol{l}} \hat{C}_{\boldsymbol{l},r,\boldsymbol{p}}\;\Upsilon_{\boldsymbol{l}}(j)
    = \prod_f \left(\sum_{l_f} \hat{C}^{(f)}_{l_f,r_f,p_f}\;\Upsilon_f[j,l_f]\right).
    \label{eq:theta_factored}
\end{equation}

The key result of \eqref{eq:theta_factored} is that $\Theta_{j,r,\boldsymbol{p}}$ is
a \emph{product} of per-factor spatial images.  Each factor image
$\varphi^{(f)}_{r_f,p_f}(j) = \sum_{l_f} \hat{C}^{(f)}_{l_f,r_f,p_f}\,\Upsilon_f[j,l_f]$
requires only $L_f$ spatial additions, so $\Theta$ can be formed without ever storing
$L_{\rm fwd} = L_o\!\times\!\prod_q L_q$ full-resolution images simultaneously.

% ---------------------------------------------------------------------------
\section{Operation Count}

Table~\ref{tab:ops} summarises the number of Fourier-domain operations per CG iteration
for the joint and factored approaches with $L_o = 7$, $L_q = 4$ for $Q=2$ NL fields.

\begin{table}[h]
\centering
\caption{NUFFTs and Toeplitz products per CG iteration (one coil, one frame).}
\label{tab:ops}
\begin{tabular}{lcc}
\toprule
Method & Forward NUFFTs & Normal-op Toeplitz prods \\
\midrule
Joint SVD (rank $L$) & $L = 112$ & $P \times R$ (large $P$) \\
Factored SVD & $L_o \!\times\! L_1 \!\times\! L_2 = 112$ & $P_{\rm total} \times R_{\rm total}$ (small) \\
Joint with $L=15$ & $15$ (poor accuracy) & $P \times R$ \\
Factored with $L_o{+}L_q{+}L_q = 15$ & 15 (same count via sharing) & $P_{\rm total} \times R_{\rm total}$ \\
\bottomrule
\end{tabular}
\end{table}

The forward-NUFFT count is identical for joint and factored at the same accuracy.
The \emph{key advantage} is:
\begin{enumerate}
\item \textbf{Accuracy at small $L$:} because the error bound is additive
(Theorem~1), factored achieves Fessler-level off-resonance accuracy with its
dedicated $L_o = 7$ vectors even when the NL correction is imperfect, whereas
joint SVD at rank 15 fails to reach off-resonance accuracy at all.

\item \textbf{Secondary compression:} the per-factor matrices in
\eqref{eq:factor_svd2} are much smaller ($L_f(L_f+1)/2 \times K$ with $L_f\leq 7$)
than the full compound matrix ($L_{\rm fwd}(L_{\rm fwd}+1)/2 \times K$), making
the secondary SVD both cheaper to compute and producing smaller $P_f$.
Empirically $P_o \sim 3$, $P_q \sim 2$, $R_f = 1$, giving
$P_{\rm total} = 3 \times 2 \times 2 = 12$ Toeplitz products.

\item \textbf{Memory:} the per-factor $\Omega_f$ and $\Upsilon_f$ tensors are
stored separately and the compound tensors are formed on-the-fly.
\end{enumerate}

% ---------------------------------------------------------------------------
\section{Sharing the Off-Resonance Factor}

In 4D-Flow with cardiac frames, respiratory phases, and multiple velocity encodings,
the off-resonance field $z(\mathbf{r})$ is the \emph{same} for all encodings.
The off-resonance SVD ($\Omega_o$, $\Upsilon_o$) is therefore computed once and
shared.  Only the NL factors $\Omega_q$ change between encodings (because the
readout-gradient direction changes), so those small per-encoding SVDs are cheap.

The total forward NUFFT count per iteration is then
$N_c \times N_{\rm frames} \times N_{\rm enc} \times L_{\rm fwd}$,
where for $N_c = 32$ coils, $N_{\rm frames} = 20 \times 8 = 160$ frames,
$N_{\rm enc} = 5$ encodings, and $L_{\rm fwd} = L_o \times L_q^2 = 7 \times 16 = 112$:
$32 \times 160 \times 5 \times 112 = 28.7\,\text{M}$ NUFFTs per iteration.
With $L_q = 2$ (valid for nl\_scale $\lesssim 0.3$, realistic clinically):
$32 \times 160 \times 5 \times (7 \times 4) = 7.2\,\text{M}$ NUFFTs per iteration.

% ---------------------------------------------------------------------------
\section{Relationship to the Existing MRI\_Physics.tex Framework}

The framework in \texttt{MRI\_Physics.tex} already covers the secondary compression
in full generality.  The only change needed is to replace the single Omega/Upsilon
pair from a joint SVD with the compound product
$(\Omega_{\boldsymbol{l}}, \Upsilon_{\boldsymbol{l}})$ from the factored SVD.
Equations~(1)--(6) of \texttt{MRI\_Physics.tex} apply verbatim with the
substitution $l \to \boldsymbol{l}$, $L \to L_{\rm fwd}$, $\Omega_{kl} \to
\Omega_{\boldsymbol{l}}(k)$, and $\Upsilon_{jl} \to \Upsilon_{\boldsymbol{l}}(j)$.

The histogram becomes \emph{per-factor}: instead of one joint histogram over
$(z, b_1, \ldots, b_Q)$, we maintain $1 + Q$ separate histograms, each of which
is one- or two-dimensional.  This also makes the phi operator construction cheaper
and produces better-conditioned SVD problems.

\end{document}
